\documentclass[10pt,a4paper]{amsart}
\usepackage[margin=19mm]{geometry}
\usepackage{amsmath,amssymb,mathtools}
\usepackage[scr=rsfs,cal=euler]{mathalfa}
\usepackage{xfrac}
\usepackage[unicode,hidelinks,bookmarks=false]{hyperref}
\newcommand{\ie}{i.e.\ }
\newcommand{\cf}{cf.\ }
\newcommand{\resp}{respectively\ }
\newcommand{\Subsection}{\subsection}
\providecommand{\nobreakdash}{\nobreak\hbox{-}}
\setlength{\emergencystretch}{2em}
\multlinegap0pt
\usepackage{amssymb}
\usepackage{xcolor}
\usepackage{float}
\newtheorem{proposition}{Proposition}
\newcommand{\B}{\mathcal{B}}
\newcommand{\C}{\mathbb{C}}
\newcommand{\F}{\mathcal{F}}
\title{On some indices of foliations and applications}
\author{Arturo Fern\'andez-P\'erez, Evelia R. Garc\'ia Barroso, Nancy Saravia-Molina}
\date{Source: arXiv:2506.01090v2 (CC BY 4.0)}
\begin{document}
\maketitle

\noindent\emph{Source and licence.} On some indices of foliations and applications. Version arXiv:2506.01090v2 (CC BY 4.0), distributed by arXiv under the Creative Commons Attribution 4.0 International licence, http://creativecommons.org/licenses/by/4.0/, as recorded in the arXiv licence metadata for this exact version and shown on the abstract page https://arxiv.org/abs/2506.01090v2. Full licence notice: \texttt{licenses/arxiv-2506.01090-CC-BY-4.0.txt}.\par\smallskip

\noindent\textit{Editorial scope.} Counted unit: the article's proposition iguales (source0.tex lines 348--355). The article's complete original proof of it is delivered with it (source0.tex lines 356--360, one \texttt{proof} environment of 5 lines). No mathematics is written, repaired or re-derived: the statement and the proof are the article's own lines, in the article's own notation, and no retained line is substituted or re-flowed.  Retained uncounted as the source-local context the proof needs: lines 238--247; lines 330--332; lines 338--340.  Nothing was omitted from any retained span, so the package carries no omission record.  No retained line contains a citation, so the wrapper carries no bibliography and no bibliography entry.  No retained line contains a figure environment, an image-inclusion command or drawing code, so no image is delivered and the package declares no asset dependency.  Editorial formatting only, in addition: an AMS \texttt{amsart} preamble that restates the article's own declarations and macros used by the retained lines, namely the environments \texttt{proposition} on separate counters, together with the standard packages those lines need.  The retained environments print in this document as Proposition 1 in order of appearance, whereas the article prints the same items with the numbering of its own full document.  The complete native source of the article as distributed in the arXiv e-print for this exact version is bundled unchanged as \texttt{sources/179136/source0.tex}.
\begin{proposition}\label{prop_1}
Let $\F$ be a germ of a singular holomorphic foliation at $(\C^2,p)$ and let $C$ be any $\F$-invariant reduced curve. Then 
\begin{eqnarray}\label{eq_5}
GSV_p(\F,C)=\mu_p(\F,C)-\mu_p(C)
\end{eqnarray}
and 
\begin{eqnarray}\label{eq_impor}
\mu_p(\F,C)-\tau_p(\F,C)=\mu_p(C)-\tau_p(C).
\end{eqnarray}
\end{proposition}

\begin{equation}\label{mudivisor}
\mu_{p}(\mathcal{F},\mathcal{B})=\mu_{p}(\mathcal{F},\mathcal{B}_0)-\mu_{p}(\mathcal{F},\mathcal{B}_{\infty})+1.
\end{equation}

    \begin{equation} \label{eq:IaMC}
    \mu_{p}(\mathcal{F},\mathcal B)=\mu_{p}(\mathcal F)-\chi_p(\mathcal{F}).
    \end{equation}

\begin{proposition}\label{iguales}
Let $\F$ be a germ of a singular foliation  at $(\C^2,p)$. Let $\B=\B_0-\B_{\infty}$ be a primitive balanced divisor of separatrices for $\F$ at $p$.  Then,
 \begin{equation*}
 \mu_p(\F)-\tau_p(\F,\B_0)=\mu_p(\B_0)-\tau_p(\B_0)-\mu_p(\F,\B_{\infty})+\chi_p(\F)+1.
 \end{equation*}
 Moreover, $\mu_p(\F)=\tau_p(\F,\B_0)$ if and only if  \[\mu_p(\B_0)-\tau_p(\B_0)=\mu_p(\F,\B_{\infty})-\chi_p(\F)-1.\]
In particular, if $\B$ is an effective divisor then $\mu_p(\F)=\tau_p(\F,\B_0)$ if and only if  $\mu_p(\B_0)=\tau_p(\B_0)$ and $\chi_p(\F)=0$.
\end{proposition}

\begin{proof}
It follows from \eqref{eq:IaMC} and \eqref{mudivisor} that $\mu_{p}(\mathcal F)=\mu_{p}(\mathcal{F},\mathcal{B}_0)-\mu_p(\F,\B_{\infty})+\chi_p(\mathcal{F})+1$. The proof follows by applying Proposition \ref{prop_1} to $\B_0$, and substituting into the last equation.
Moreover, $\mu_p(\F)=\tau_p(\F,\B_0)$ if and only if \[\mu_p(\B_0)-\tau_p(\B_0)=\mu_p(\F,\B_{\infty})-\chi_p(\F)-1.\]
If $\B$ is effective, then the  last part of the proposition follows since $\mu_p(\B_0)-\tau_p(\B_0)$ and $\chi_p(\F)$ are non-negative integers.
\end{proof}

\end{document}
